% !TEX root = ../main.tex

\documentclass[../main.tex]{subfiles}

\graphicspath{{\subfix{../images/}}}

\begin{document}

\subsection{Используемые библиотеки}

\begin{table}[H]
\centering
\caption{Внешние зависимости проекта}
\label{tab:deps}
\begin{tabularx}{\textwidth}{|l|l|X|}
\hline
\textbf{Библиотека} & \textbf{Версия} & \textbf{Роль в проекте} \\
\hline
MAVLink \texttt{c\_library\_v2}~\cite{mavlink_docs} & 2.0
    & Парсинг телеметрии, декодирование \texttt{RAW\_IMU} \\
\hline
GLFW3~\cite{glfw_docs} & $\geq$3.3
    & Создание окна OpenGL, ввод мыши и клавиатуры \\
\hline
GLEW & $\geq$2.1
    & Загрузка расширений OpenGL~3.3 \\
\hline
GLM  & $\geq$0.9.9
    & Матричная и кватернионная математика \\
\hline
pthreads & POSIX
    & Межпоточная синхронизация (\texttt{std::thread}, \texttt{std::mutex}) \\
\hline
\end{tabularx}
\end{table}

Библиотека MAVLink является заголовочной~--- она полностью состоит из
\texttt{.h}-файлов и не требует компиляции~\cite{mavlink_proto}.
Остальные зависимости установливаются из apt-репозиториев.

\subsection{Конфигурация CMake}

Сборочная система использует CMake версии не ниже~3.10 и настроена на
стандарт C++14. Зависимости ищутся через \texttt{find\_package} и
\texttt{PkgConfig}; проект компилируется с флагами \texttt{-O2 -Wall -Wextra}.
CMake ищет GLFW тремя способами последовательно: через \texttt{find\_package},
через \texttt{PkgConfig} и, наконец, вручную через \texttt{find\_path} /
\texttt{find\_library}~--- это обеспечивает совместимость с различными
дистрибутивами GNU Linux.

	
\end{document}
